\makeatletter \def\input@path{{../../}} \makeatother
\documentclass{article}
\usepackage[utf8]{inputenc}
\usepackage{amsmath, amsthm, amssymb, mathpazo, isomath, mathtools}
\usepackage{subcaption,graphicx,pgfplots}
\usepackage{fullpage}
\usepackage{booktabs}
\usepackage{hyperref}
\usepackage{algorithm, algorithmic}
\usepackage{mathtools}
\usepackage{todonotes}
\usepackage{tcolorbox}
\usepackage{totcount,xfp}
\newtotcounter{totalpts}
%\setcounter{totalpts}{0} % Contador de puntos en homework

\newcommand{\getTotal}{\number\totvalue{totalpts} }
\newcommand{\instruccion}{\begin{tcolorbox}[colback=blue!5!white,colframe=blue!75!black!70!white,title=Instructions]
The homework must be developed by one or two people in a .pdf report submitted through the designated drop box on the Canvas platform, where you must also show the developed code if applicable. For your convenience, you may submit the homeworks in a Jupyter Notebook, so it is more comfortable to show the code. In any case, a single file must be submitted as the answer to the homework. The late policy will be: a linear factor will be calculated that is 1 at the time of delivery and 0 48 hours later. This will multiply your obtained score.


The homework has \getTotal points. In each question, the given score will be: the total if it is correct, half if it has a conclusive development but contains errors, and 0.0 points if it has little to no development.

The use of language models (like ChatGPT) is allowed to support the development of the homework, although we suggest trying each question a bit first to generate doubts. While we cannot control if you ask the model directly for the answer, we believe it is better for your learning process to ask for suggestions rather than solutions. For this reason, if you decide to use it, you must attach to your homework another pdf showing the prompts used. We suggest using prompts of this style: "I am a student of [Major] and I am developing the following exercise for a homework in the course [Course]: [Copy exercise]
  And we have seen so far the contents in the attached slides [attach pdfs of the lectures]. Given this content, I would like to ask you for a suggestion to solve [explain difficulty], since I am confused about [write doubt]. As an answer, I do not want you to give me the solution, but a guide to help me start developing my answer."

  If you have doubts and/or comments, please contact us by email or Canvas. It is possible that the formulation has errors.
\end{tcolorbox}}


\newcommand{\code}[1]{\texttt{#1}}
% Logic: To create a counter with the points in each question and accumulate that in a total variable. Now I simply set the full threshold at the beginning to 75% and it gives the computed value automatically.
\newcommand{\pts}[1]{\if#11\textbf{[#1 point]}\else\textbf{[#1 points]}\fi\addtocounter{totalpts}{#1}}


\title{Homework 5}
\author{}
\date{IMT3130}

\renewcommand{\vec}{\vectorsym}
\newcommand{\mat}{\matrixsym}
\newcommand{\ten}{\tensorsym}
\DeclareMathOperator{\grad}{\nabla}
\DeclareMathOperator{\dive}{\text{div}}
\DeclareMathOperator{\curl}{\text{curl}}
\DeclareMathOperator{\tr}{\text{tr}}
\DeclareMathOperator{\sym}{\text{sym}}
\newcommand{\R}{\mathbb{R}}
\newcommand{\D}{\mathcal{D}}

\begin{document}

\maketitle

\instruccion

\begin{enumerate}
    \item Consider the deformation given by 
        \begin{align*}
            x_1(\vec X, t) &= e^tX_1 - e^{-t} X_2 \\
            x_2(\vec X, t) &= e^t X_1 + e^{-t}X_2 \\
            x_3(\vec X, t) &= X_3. 
        \end{align*}
    \begin{itemize}
        \item\pts{1} Compute the displacement field, the velocity, and the acceleration in material (reference) coordinates. 
        \item\pts{1} Compute the displacement field, the velocity, and the acceleration in spatial (deformed) coordinates. 
        \item\pts{1} Compute the tensors $\ten F$, $\ten C$, and $\ten E$.
        \item\pts{1} Ignore the third component and consider $\Omega$ a unit square. Plot in Python the initial configuration ($t=0$), and the deformed configuration at times 0.5, 1.0, and 2.0. 
        \item\pts{1} At the described times, compute the Frobenius norm of $\ten C$ in the reference configuration and in the deformed one. Comment on how the interpretation of the tensor changes between both configurations, and propose a way to modify the tensor $\ten C$ so that it has a correct interpretation in the deformed configuration.
    \end{itemize}

    \item Consider a Neo-Hookean hyperelastic material, where the hyperelastic energy is given by
            $$ \Psi(\ten F) = \frac C 2(\tr \ten C - 3) + \frac \kappa 2 \left( J -1\right)^2. $$
            \begin{itemize}
                \item\pts{1} Compute the first and second Piola-Kirchhoff stress tensors.
                \item\pts{1} Write the weak formulation associated with the momentum equation with the computed explicit Piola tensor and appropriate boundary conditions. Since it is a non-linear problem, it is not clear a-priori what the correct regularity space for this model is. To develop the exercise, simply consider $H^1$ as the ambient space.
            \end{itemize}
    \item Modify the derivation of mass conservation to consider a mass flux given by its normal gradient, i.e. consider an additional surface term for any subdomain $\omega_t\subseteq\Omega_t$: 
            $$ \int_{\partial \omega_t} \ten K \grad\rho \cdot \vec n \,dS, $$
            \begin{itemize}
                \item\pts{1} Show the resulting equation, and show that the spatial differential operator is an ADR operator. 
            \end{itemize}
    \item Consider the reference linear momentum equation: 
            $$ \rho_0\ddot{\vec u} - \dive_X \ten P(\ten F)=0. $$
            \begin{itemize}
                \item\pts{1} Rewrite the tensor $\ten P$ so that it depends on $\ten E$, and then linearize with respect to $\vec u = 0$. Remember that linearizing an equation
                    $$ F(x) = 0 $$
                    with respect to a point $x_0$ means replacing $F(x) \approx F(x_0) + dF(x_0)[\delta x]$ to obtain an equation for the increment $\delta x$. Here, $dF(x_0)[\delta x]$ is the Gâteaux derivative in the direction $\delta x$. 
                \item\pts{1} Show that the resulting problem is the linear elastodynamics problem, and clearly identify the obtained Hooke tensor. 
            \end{itemize}

    \item Consider the incompressibility constraint $J=1$ for non-linear elastic problems. 
        \begin{itemize}
            \item\pts{1} Write the variational principle of hyperelasticity restricted to $J=1$, and then show the saddle-point problem with a Lagrange multiplier that helps impose the condition $J=1$ through an unconstrained problem.
            \item\pts{1} Show the first-order conditions for the two-variable system obtained in the previous point.
        \end{itemize}


    \item\pts{1} Define your course project. To do this, consider the following breakdown: 
            \begin{enumerate}
                \item Delimit the problem of interest and explain its origin. Deriving a model can also be part of the project, as long as what will be done is well explained. 
                \item The project must contemplate a discretization process.
                \item The project must contemplate the stability and convergence analysis of the proposed discretization.
            \end{enumerate}

            Like the homeworks, the projects are in groups of up to 2 people.


\end{enumerate} % preguntas

\end{document}
